% concatenated from On_quandle_representations.tex
\documentclass[12pt, reqno]{amsart}


\usepackage[utf8]{inputenc}
\usepackage[T1]{fontenc}
\usepackage{tikz,tikz-3dplot}
\usepackage{amsmath,amsfonts,amssymb,amsmath,latexsym,mathrsfs}
\usepackage{shuffle}
\usepackage[all]{xy}
\usepackage{stackengine}
 
\usepackage{enumerate}

\usepackage{hyperref}
 
\usepackage{tabu}
\usepackage{tikz-cd} 

\usepackage{comment}
\usepackage{url}
\usepackage{tikz}

\usepackage{xfrac}    
 
 
\setlength{\oddsidemargin}{0in}

\setlength{\evensidemargin}{0in}

\setlength{\textheight}{9in}

\setlength{\textwidth}{6.5in}

\setlength{\topmargin}{-0.5in}

\pagestyle{plain}

\newtheorem{theorem}{Theorem}[section]

\newtheorem{proposition}[theorem]{Proposition}
\newtheorem{definition}[theorem]{Definition}

\newtheorem{corollary}[theorem]{Corollary}

\newtheorem{remark}[theorem]{Remark}
\newtheorem{conjecture}[theorem]{Conjecture}

\newtheorem{lemma}[theorem]{Lemma} 

\newtheorem{example}[theorem]{Example} 
\newcommand{\N}{ {\mathbb N} }
\newcommand{\Z}{ {\mathbb Z} }
\newcommand{\C}{ {\mathbb C} }
\newcommand{\K}{ {\Bbbk} }
\newcommand{\e}{ {\epsilon} }
 \newcommand{\Q}{{\mathbf{Qu}}}
 \newcommand{\G}{{\mathbf{Grp}}}
\begin{document}

 
\title{On quandle representations}
\keywords{Quandles, enveloping groups, representations}

 


 
\author{Mohamad \textsc{Maassarani} }
\maketitle
\begin{abstract} A unitary finite dimensional quandle representation is decomposable into a direct sum of irreducible represenations. Not all quandle representations satisfy this property. We prove that a finite dimensional quandle represenation $\rho :Q \to GL(V) $ of a finite quandle $Q$ over $\mathbb{C}$ is decomposable into a direct sum of irreducibles if and only if every element in the image of $\rho$ is diagonlizable. We show that an irreducible representation $\rho :Q \to GL(V)$ of a finite quandle over $\mathbb{C}$ is unitary for some inner product if and only if every element of the image of $\rho$ has determinant of modulus $1$. It follows that any irreducible representation of a finite quandle $Q$ over $\mathbb{C}$ can be twisted by a quandle character to obtain a unitary irreducible representation. We also prove that the enveloping group $G(Q)$, of a finite quandle $Q$, admit a faithfull finite dimensional unitary representation over $\mathbb{C}$ and that the irreducible representations of a finite quandle $Q$ over $\mathbb{C}$ are $1$-dimensional if and only if $G(Q)$ is abelian. Finaly, we determine the irreducible representations over $\mathbb{C}$ of a family of finite quandles.
\end{abstract}
\section*{Introduction and main results}
A quandle is a set $Q$ equipped with a binary operation $\triangleright: Q\times Q \to Q$ satisfying some axioms. A group $G$ equipped with the binary operation $x\triangleright y=xyx^{-1}$ defines a quandle called the conjugacy quandle of $G$. A quandle representation of a quandle $Q$ over a vector space $V$ is a quandle morphism from $Q$ to the conjugacy quandle of $GL(V)$, i.e. $\rho(x\triangleright y)=\rho(x)\rho(y)\rho(x)^{-1}$. A subrepresentation of $\rho$ is a subspace of $V$ stable under the image of $\rho$ and $\rho$ (resp. a subrepresentation) is called irreducible if the only subrepresentations are $0$ and $V$ (respectively the subrepresentation). A quandle character of $Q$ is a quandle morphism $\chi : Q\to \mathbb{C}^\times$ where $\mathbb{C}^\times$ is considered as a conjugacy quandle, i.e. $\chi(x\triangleright y)=\chi (y)$. A quandle representation is unitary if its image preserves an inner product. 
\subsection*{Outline of the paper}These notes are divided into $5$ sections.\\\\
In section $1$, we review definitions and properties of quandles, quandle morphisms, conjugacy quandle, inner automorphism groups of quandles, quandle representations and the enveloping group of a quandle.\\\\
In section $2$, we prove that a finite dimensional unitary quandle representation is decomposable into a direct sum of irreducible subrepresentations and that a finite dimensional quandle representation $\rho$ of a finite quandle $Q$ over $\mathbb{C}$ is decomposable into direct sum of irreducibles if and only if the elements of $\rho(Q)$ are diagonalizable. \\\\
In section $3$, we prove that a finite dimensional irreducible quandle representation over $\mathbb{C}$ of a finite quandle is unitary for some inner product if and only if the determinant of any element in the image of the representation has modulus $1$. We also show that up to multiplication by a quandle character a finite dimensional irreducible representation of a finite quandle over $\mathbb{C}$ is unitary.\\\\
In section $4$, we prove that the enveloping group of a finite quandle admits a faithfull finite dimensional unitary representation over $\mathbb{C}$ and that the irreducible representations of a finite quandle over $\mathbb{C}$ are $1$-dimensional if and only if the enveloping group is abelian.\\\\
In section $5$, we determine the irreducible representations over $\mathbb{C}$ of the family of quandles $Q_{n,m}$  ($n,m\geq 1$) consisting of the set $$\{x_i \vert i \in \mathbb{Z}/n\mathbb{Z} \} \cup \{y_i \vert i \in \mathbb{Z}/m\mathbb{Z} \}$$ equipped with the binary operation defined by :
$$x_i \triangleright y_j=y_{j+1},\quad y_i\triangleright x_j=x_{j+1} ,\quad x_{i}\triangleright x_j=x_j,\quad y_i\triangleright y_j=y_j.$$

 
\section{Reminders on quandles}
\subsection{Quandle}
A quandle is a set $Q$ equipped with a binary operation $\triangleright$  such that :
\begin{itemize}
\item[-] $x\triangleright x=x$ for $x\in Q$.
\item[-] For all $x,y \in Q$ there exist a unique $z\in Q$, such that $x \triangleright z=y$.
\item[-] $x \triangleright( y\triangleright z)= (x \triangleright y) \triangleright (x \triangleright z)$, for $x,y,z \in Q$. 
\end{itemize}
The conjugacy quandle $Conj(G)$ of a group $G$ is the set $G$ equipped with the binary operation defined by $x\triangleright y= xyx^{-1}$  for $x,y\in G$.\\\\  
A map between two quandles $(Q_1,\triangleright_1)$ and $(Q_2 ,\triangleright_2)$ satisfying : $$f(z \triangleright_1 w)=f(z) \triangleright_2 f(w),$$
for $z,w \in Q$ is called a quandle mophism. Composition of quandle morphisms gives quandle morphisms. A group morhism is a quandle morphism if the groups are considered as conjugacy quandles.
\subsection{Inner automorphism group}
 For $Q$ a quandle and $x\in Q$, the left translation $L_x$ given by $L_x(y)=x\triangleright y$ for $y\in Q$ is a bijective quandle morphism. The inner automorphism group $Inn(Q)$ of a quandle $Q$ is the subgroup of bijections of $Q$ generated by left translations. The map from $Q$ to $Inn(Q)$ assigning to an element the corresponding left translation is quandle morphism with respect to the conjugacy quandle structure on $Inn(Q)$.
\subsection{Quandle representation}
Quandle representations where introduced in \cite{rep}. A representation of a quandle $Q$ is a vector space $V$ and a quandle morphism $\rho: Q \to Conj(\mathrm{GL}(V))$, i.e. : 
$$\rho(x\triangleright y)=\rho(x) \rho(y) \rho(x)^{-1},$$
for all $x,y \in Q$.
A subrepresentation of $\rho$ is a vector subspace $W\subset V$, stable under the elements of $\rho(Q)$.  An irreducible representation is a representation $V$ that has no subrepresentations other than $0$ and $V$. We define intertwining operators and equivalent representations as for group representations.\\\\
In \cite{Qrep}, we define a quandle character of a quandle $Q$ as a quandle morphism $\chi : Q \to \mathbb{C}^\times$, i.e. :$$\chi(x\triangleright y) =\chi(y),$$ for $x,y\in Q$. Quandle characters of $Q$ correspond to functions of $Q$ into $\mathbb{C}^\times$ that are constant on the orbits of $Q$ under $Inn(Q)$.\\\\
A quandle representation $\rho : Q \to GL(V)$ will be called unitary if $V$ is endowed with an inner product for wich the elements of $\rho(Q)$ are isometries.
\subsection{Enveloping group}
The enveloping group $G(Q)$ of a quandle $Q$ is the group :
$$G(Q)=\langle x \in Q \vert xyx^{-1}=x\triangleright y , \text{ for } x,y\in Q\rangle.$$
The enveloping group is also called in the litterature : the associated group, adjoint group or structure group. The group $G(Q)$ is infinite and the map $\varphi_Q : Q \to G(Q)$ assigning to $x\in Q$ the corresponding generator is a quandle morphism for the conjugcay quandle structure on $G(Q)$. $\varphi_Q$ is universal with respect to quandle morphisms into conjugacy quandles : for $f:Q\to Conj(G)$ a quandle morphism there is a unique group morphism $\tilde{f}:G(Q)\to G$ such that $f=\tilde{f}\circ\varphi_Q$. \\\\
If $f:Q\to Q'$ is a quandle morphism between two quandles there is a unique group morphism $G(f):G(Q)\to G(Q')$ such that the following diagram commutes :
$$\begin{tikzcd}
G(Q) \arrow{r}{G(f)} &G(Q')  \\ 
Q\arrow{u}{\varphi_Q} \arrow{r}{f} & Q'\arrow{u}{\varphi_{Q'}} 
\end{tikzcd} $$
The assignement $Q\mapsto G(Q)$, $f\mapsto G(f)$ is functorial.
 

\section{Decomposition of quandle representations}
\begin{definition}
A quandle representation will be called completely reducible if it can be decomposed as the direct sum of irreducible subrepresentations.
\end{definition}

\begin{theorem}
A finite dimensional unitary quandle representation is completely reducible.
\end{theorem}
\begin{proof}
Let $V$ be the vector space of the representation and $Q$ the quandle. For dimensional reasons $V$ contains an irreducible subrepresentations $V_1$. Since $V_1$ is stable under the elements of $Q$ the orthogonal $V^1$ of $V_1$ is stable under the elements of $Q$. $V$ hence decomposes as the sum of two subrepresentations $V_1\oplus V^1$ with $V_1$ irreducible and $V^1$ unitary. For dimensional reasons $V^1$ has an irreducible subrepresentation $V_2$ and the orthogonal $V^2$ of $V_2$ in $V^1$ is stable under the element of $Q$. We now have $V=V_1\oplus V_2 \oplus V^2$. Iterating the process we get that $V$ decomposes as a direct sum of irreducible subrepresentations.
\end{proof}
\begin{example}
Let $Q$ be a finite quandle and $R\subset Q$ such that for $x\in Q$ and $y\in R$, $x\triangleright y \in R$ $($i.e. $R$ is a union of orbits under the action of $Inn(Q)$$)$. For $x\in Q$ the left translation $L_x$ induces a permutation of $R$. Hence, it induces a unitary automorphism $\rho_R(x)$ of the free complex vector space $\mathbb{C}R$ over $R$. Since $L_{x\triangleright x'}=L_xL_{x'}L_x^{-1}$. The assignement $\rho_R :Q \to GL(\mathbb{C}R), x\mapsto\rho_R(x)$ is a unitary quandle representation. For $R=Q$ the representation $\rho_R$ and the regular representation of $Q$ introduced in \cite{rep} are dual.
\end{example}
As noticed in \cite{nondec}, in general a quandle representation is not always completely reducible. Indeed, if $Q$ is a quandle the mapping $\rho :Q \to GL(\mathbb{C}^2)$ wich assigns to every element of $Q$ the endomorphism having the following matrice in the canonical basis :
\[\begin{pmatrix}
1 & 1\\
0 & 1
\end{pmatrix}\]
is a quandle representation. As one checks this representation is not completely reducible.\\\\
Let $\rho : Q \to GL(V)$ be a quandle representation, by the universal property of the map $\varphi_Q$ from $Q$ to the envelopping group $G(Q)$ there is a unique group representation $\rho_{G(Q)}: G(Q)\to GL(V)$ such that $\rho_{G(Q)}=\rho \circ \varphi_Q$.
\begin{proposition}
\item[1)] The assignement $\rho \mapsto \rho_{G(Q)}$ is a $1$ to $1$ correspondance between quandle representations of $Q$ and group representations of $G(Q)$.
\item[2)] A subrepresentation for $Q$ with respect to $\rho$ is a subrepresentations for $G(Q)$ with respect to $\rho_{G(Q)}$ and vice versa.
\end{proposition}
\begin{corollary}\label{corred}
The quandle representation $\rho :Q\to GL(V)$ is completely reducible if and only if the corresponding group representation $\rho_{G(Q)} :G(Q) \to GL(V)$ is completely reducible.
\end{corollary}
\begin{corollary}
$\rho$ is irreducible if and only if $\rho_{G(Q)}$ is.
\end{corollary}


\begin{proposition} Let $Q$ be a finite quandle and denote by $Q/Inn(Q)$ the set of orbits of $Q$ with respect to the action of $Inn(Q)$.
\item[1)] The abelianisation of the enveloping group $G(Q)$ is isomorphic to $\mathbb{Z} Q/Inn(Q)$, the free abelian group on $Q/Inn(Q)$.
\item[2)] The abelianisation map with the idetification in $1)$ $Ab :G(Q)\to \mathbb{Z} Q/Inn(Q)$ maps a generator to its orbit in $Q/Inn(Q)$.
\item[3)] The derived group of $G(Q)$ is finite. 
\end{proposition}
\begin{proof}
$1)$ and $2)$ are well established results. For instance one can prove them using generators and relations. Now, the group $G(Q)$ is a central extension of $Inn(Q)$ (\cite{EM}). Since $Q$ is finite $Inn(Q)$ is finite and the center of $G(Q)$ has finite index. $3)$ is hence an application of Schur's theorem for the derived group.
\end{proof}

\begin{proposition}\label{Z0} For $Q$ a finite quandle let $Z_0$ be the subgroup of $G(Q)$ generated by the elements $x^n$ for $x\in Q$ and with $n=\vert Inn(Q) \vert $.
\item[1)] The group $Z_0$ lies in the center of $G(Q)$
\item[2)] $Z_0$ has finite index in $G(Q)$.
\end{proposition}
\begin{proof}
For $x$ in $Q$, we have denoted, in the first section, by $L_x$ the left translation by $x$ (bijection of $Q$ mapping $y$ to $x\triangleright y$). From what we have seen in the first section, there exists a unique group morphism $G(L_x):G(Q)\to G(Q)$ such that $G(L_x)\circ \varphi_Q=\varphi_Q \circ L_x$. As one can check on the generators of $G(Q)$, the conjugacy $c_x$ in $G(Q)$ by the generator $x$ satisfies the last condition. Hence, $G(L_x)=c_x$. Since the assignement is functorial, $G(L_x^n)=G(id_Q)=c_x^n=id_{G(Q)}$ where $n$ is as in the proposition. This proves that for $x\in Q$, $x^n$ lies in the center of $G(Q)$. We have proved $1)$. It can be readly checked from $1)$ and $2)$ of the previous propostion that the image of $Z_0$ in the abelianisation of $G(Q)$ is of finite index. By $3)$ of the previous proposition the derived group of $G(Q)$ is finite and therefore $Z_0$ has finite index in $G(Q)$.
\end{proof}

\begin{lemma}\label{power}
An endomorphism $f$ of a finite dimensional vector space over $\mathbb{C}$ is diagonalizable if and only if $f^n$ is diagonalizable for some $n>0$.
\end{lemma}
\begin{proof}
If $f$ is diagonalizable then $f^n$ is diagonalizable. If $f^n$ is diagonalizable the minimal polynomial $P(X)$ of $f^n$ has simple roots and so does the polynomial $R(X)=P(X^n)$. But, $R(f)=P(f^n)=0$ and $R$ has simple roots, hence $f$ is diagonalizable. This proves that if $f^n$ diagonalizable then $f$ is.
\end{proof}

\begin{proposition}\label{dec}
Let $1\to Z \to G \to H\to 1$ be a central extension of groups with $H$ finite and $\rho :G\to GL(V)$ be a finite dimensional complex group representation.
\item[1)] The elements of $\rho(Z)$ are all diagonalizable if and only if all elements of $\rho(G)$ are diagonalizable.
\item[2)] The representation $\rho$ is completely reducible if and only if every element of $\rho(G)$ $($equivalently $\rho(Z)$$)$ is diagonalizable.
\end{proposition}
\begin{proof}
We prove $1)$. For $g\in G$, the element  $g^{\vert H \vert }$ lies in $Z$ and hence by the previous lemma $\rho(g)^{\vert H \vert }$ is diagonalizable if and only if $\rho(g)$ is. $1)$ follows.
We now prove $2)$. If $\rho$ is decomposable into sum of irreducibles $\bigoplus_i V_i$ then by Schur's lemma $Z$ acts by multiplication by scalar in each $V_i$ and hence any element of $\rho(Z)$ is diagonalizable. We have proved that if $\rho$ is decomposable into sum of irreducibles then $\rho(z)$ is diagonalizable for every $z\in Z$. We now prove the converse. Assume that all the elments $\rho(z)$ for $z\in Z$ are diagonalizable. This means that the elements of $\rho(Z)$ can be simultaneously diagonalized since they commute. Hence, there is a basis $(v_1,\dots,v_n)$ of $V$ such that for all $z\in Z$ : 
$$ \rho(z) v_i =\chi_i(z) v_i,$$ 
with $\chi_i :Z \to \mathbb{C}^\times$ a multiplicative character. Let $X$ be the set formed by the characters $\chi_1,\cdots,\chi_n$. For $\chi \in X$, let $V_\chi$ be the subspace of $V$ spanned by the vector $v_i$ such that $\chi_i=\chi$. The space $V_\chi$ is stable under the action of $G$ and :
$$ V= \bigoplus_{\chi\in X} V_{\chi}.$$
For $\chi \in X$, denote by $\rho_\chi : G\to GL(V_\chi)$ the representation obtained from $\rho$ by restriction to $V_\chi$ and by $\varphi_\chi$ the projection $GL(V_\chi)\to PGL(V_\chi)$. The subgroup $Z$ of $G$ acts in $V_X$ by multiplication by scalars and $G$ is the central extension of the finite group $H$ by $Z$. Hence, the image of $\varphi_\chi \circ \rho_\chi$ is finite. Since we are working over $\mathbb{C}$, the projection $\theta_\chi : SL(V_\chi)\to PGL(V_\chi)$ ($SL$ for endomorphisms of determinant $1$) is surjective and has finite kernel. Hence, the preimage $G_\chi$ of $\varphi_\chi \circ \rho_\chi(G)$ under $\theta_\chi$ is a finite group. Since $G_\chi$ is finite, $V_\chi$ decomposes as sum of irreducibles under the natrual action of $G_\chi \subset SL(V_\chi)$ : 
$$V_\chi = \bigoplus_i W_i$$
 Now any element of $\rho_\chi(G)$ is obtained from an element of $G_\chi$ by a multiplication by a scalar and vice versa. Hence, any stable subspace of $V_\chi$ under $G_\chi$ is stable under $G$ and vice versa. Therefore, the decomposition of $V_\chi= \bigoplus_i W_i$ abtained above is a decomposition of $V_\chi$ into irreducibles with respect to the representation $\rho_\chi :G \to GL(V_\chi)$. This proves that $V$ decomposes into a direct sum of irreducible representations if $\rho(z)$ is diagonalizable for all $z\in Z$. We have completed the proof of $2)$ of the proposition.
\end{proof}

\begin{theorem}
Let $Q$ be a finite quandle and $\rho :Q \to GL(V)$ be a finite dimensional quandle representation over $\mathbb{C}$. The quandle representation $\rho$ is completely reducible if and only if all the elements of $\rho(Q)$ are diagonalizable.
\end{theorem}
\begin{proof}
By corollary \ref{corred}, $\rho$ is completely reducible if and only if the correspondant group representation  $\rho_{G(Q)}:G(Q)\to GL(V)$ is. From proposition \ref{Z0}, the subgroup $Z_0$ of $G(Q)$ generated by $x^n$ for $x\in Q$ and $n=\vert Inn(Q)\vert$ lies in the center of $G(Q)$ and has finite index. Applying $2)$ of the previous proposition to the central exact sequence $1\to Z_0\to G(Q) \to G(Q)/Z_0 \to 1$, we get that $\rho_{G(Q)}$ is completely reducible if and only if the elements of $\rho_{G(Q)}(Z_0)$ are diagonalizable. But $Z_0$ is abelian and the diagonalizability of the image of its elements is equivalent to the diagonalizability of the image of its generators. Hence, $\rho_{G(Q)}$ is completely reducible if and only if the element $\rho_{G(Q)}(x)^n$ for $x\in Q$ with $n=\vert Inn(Q) \vert $ are diagonalizable. By lemma \ref{power}, $\rho_{G(Q)}(x)^n$ is diagonalizable if and only if $\rho_{G(Q)}(x)=\rho(x)$ is. This proves the theorem.
\end{proof}
\begin{corollary}
For $Q$ a finite quandle and $\rho : Q \to GL(V)$ a finite dimensional irreducible quandle representation over $\mathbb{C}$, the image of $\rho$ consists of diagonalizable automorphisms.
\end{corollary}
\begin{remark}
The complete reduciblity of $\rho$ as in the theorem implies the diagonalizability of all the elements in the image of $\rho_{G(Q)}$ $($by following the proof and the previous proposition$)$.
\end{remark}
\begin{example}
If $Q$ is a finite quandle, $G$ is a finite group, $f:Q\to Conj(G)$ is a quandle morphism and $\rho :G\to GL(V)$ a finite dimensional group representation of the group $G$ over $\mathbb{C}$, the map $\rho \circ f$ is a quandle representation that is completely reducible. Indeed, since $G$ is finite, the image of $\rho$ consists of diagonalizable automorphisms.
\end{example}
We note that the last theorem could fail if the quandle is not finite. Indeed, let $G$ be the subgroup of $GL_2(\mathbb{C})$ generated by the matrices :
\[\begin{pmatrix}
1 & 1\\
0 & 1
\end{pmatrix}\quad \text{and} \quad \begin{pmatrix}
1 & 0\\
1& 1
\end{pmatrix}.\]
$G$ is infinite since the matrices don't have a finite order. $Conj(G)$ is a quandle and the natural inclusion of $Conj(G)=G$ into $GL_2(\mathbb{C})$ is an irreducible quandle representation. Indeed, the only nonzero stable subspace shared by the two matrices is $\mathbb{C}^2$. But the matrices given are not diagonalizable. 
\section{Irreducible unitary representations}
\begin{theorem}
Let $Q$ be a finite quandle and $\rho : Q \to GL(V)$ be an irreducible finite dimensional complex representation of $Q$. $\rho$ is unitary for some inner product if and only if the determinant of every element of $\rho(Q)$ has modulus $1$.
\end{theorem}
\begin{proof}
If $\rho$ is unitary the elements of $\rho(Q)$ are isometries and hence have determinant with modulus $1$. We will prove the converse. Assume that the determinant of any element of $\rho(Q)$ has modulus $1$. Let $\rho_{G(Q)} : G(Q) \to GL(V)$ be the corresponding representation of $G(Q)$. Since the image of $\rho_{G(Q)}$ consist of products of elements of the image of $\rho$, the determinant of any element in the image of $\rho_{G(Q)}$ has modulus $1$. By proposition \ref{Z0}, we have a central extension $1 \to Z_0 \to G(Q) \to G(Q)/Z_0\to 1$ with $G(Q)/Z_0$ finite. Set $H=G(Q)/Z_0$. Let $s:H\to G(Q)$ be a section of $G(Q)\to H$ and $\langle-,-\rangle$ be an inner product of $V$. We define the inner product $\langle -,-\rangle_s$ on $V$ by :
$$ \langle v,w\rangle_s= \sum_{h\in H} \langle \rho_{G(Q)}(s(h))(v), \rho_{G(Q)}(s(h))(w)\rangle,$$
for $v,w\in V$. Take $g \in G$; $g$ can be written as $zs(h_1)$ for some $z\in Z_0$ and $h_1\in H$. We have that :
$$ \langle \rho_{G(Q)}(g)(v),\rho_{G(Q)}(g)(w)\rangle_s= \sum_{h\in H} \langle \rho_{G(Q)}(z_hs(hh_1))(v), \rho_{G(Q)}(z_hs(hh_1))(w)\rangle$$
where 
$$ z_h=zs(h)s(h_1)s(hh_1)^{-1},$$
lies in $Z_0$. But $Z_0$ lies in the center and the representation $\rho_{G(Q)}$ is irreducible with all elements in the image having a determinant of modulus $1$. Hence, $\rho_{G(Q)}(z_h)=\lambda_h id$ with $\lambda_h$ of modulus $1$. This proves that  :
\begin{equation*}
 \begin{split}
\langle \rho_{G(Q)}(g)(v),\rho_{G(Q)}(g)(w)\rangle_s  & = \sum_{h\in H} \langle \lambda_h\rho_{G(Q)}(s(hh_1))(v), \lambda_h\rho_{G(Q)}(s(hh_1))(w)\rangle\\
&=\sum_{h\in H}\lambda_h\overline{\lambda_h} \langle \rho_{G(Q)}(s(hh_1))(v), \lambda_h\rho_{G(Q)}(s(hh_1))(w)\rangle\\
&=\sum_{h\in H} \langle \rho_{G(Q)}(s(hh_1))(v), \rho_{G(Q)}(s(hh_1))(w)\rangle\\
&= \langle v,w\rangle_s
\end{split}
\end{equation*}
Therefore, the group representation $\rho_{G(Q)}$ is unitary with respect to $\langle -,-\rangle_s$ and so is $\rho$. We have proved the theorem.
\end{proof}
For a complex number $z=re^{i\theta}$ with $r>0$, $ 0\leq \theta < 2\pi$ and $n$ a positive natural number, we define $z^{\frac{1}{n}}$ by :
$$z^{\frac{1}{n}}=r^{\frac{1}{n}}e^{i\frac{\theta}{n}}.$$
\begin{proposition}
Let $\rho : Q \to GL(V)$ be a quandle representation with $V$ finite dimensional over $\mathbb{C}$. The map $\chi_\rho : Q \to \mathbb{C}^\times$ given by :
$$\chi_\rho(x)=\frac{1}{det(\rho(x))^{\frac{1}{dim(V)}}},$$
for $x\in Q$, is a quandle character.
\end{proposition}
\begin{proof}
This can be readly checked
\end{proof}
\begin{proposition}
Let $\rho :Q \to GL(V)$ be an irreducible finite dimensional quandle representation over $\mathbb{C}$. The representation $\chi_\rho \cdot \rho : Q  \to GL(V)$ defined by :
$$(\chi_\rho \cdot \rho)(x)=\chi_\rho(x)  \rho(x),$$
for $x \in Q$, is an irreducible unitary representation with respect to some inner product.
\end{proposition}
\begin{proof}
It is not hard to check that $\chi_\rho\cdot \rho $ is an irreducible quandle representation. The fact that it is unitary follows from the previous theorem. Indeed, the determinant of the elements of the image of $\chi_\rho \cdot \rho$ is $1$.
\end{proof}
\section{Faithfull unitary representations of enveloping groups}
\begin{lemma}
A finitely generated abelian group admits a faithfull finite dimensional unitary representation over $\mathbb{C}$.
\end{lemma}
\begin{proof}
A finitely generated abelian group is a product of cyclic groups. A cyclic group admits a faithfull $1$-dimensional unitary representation. One can construct a "product representation" from individual representations of the cyclic groups.
\end{proof}
\begin{theorem}
Let $Q$ be a finite quandle. The enveloping group $G(Q)$ of $Q$ admits a faithfull finite dimensional unitary representation over $\mathbb{C}$.
\end{theorem}
\begin{proof}
By proposition \ref{Z0}, $Z_0$ is a finitely generated subgroup of the center of $G(Q)$ moreover $Z_0$ has finite index in $G(Q)$. by the previous lemma $Z_0$ admits a faithfull finite dimensional unitary representation. The induced representation to $G(Q)$ of such representation is also faithfull, unitary and finite dimensional.  
\end{proof}
\begin{corollary}
A finite quandle such that the universal map $\varphi_Q :Q\to G(Q)$ is injective admits a faithfull $($injective$)$ unitary finite dimensional quandle representation over $\mathbb{C}$.
\end{corollary}
\begin{remark}
If $\varphi_Q$ is not injective, then $Q$ admit no faithfull representation. Indeed, a representation of $Q$ "factors" throught $G(Q)$.
\end{remark}
\begin{proposition}
Let $Q$ be a finite quandle, the irreducible quandle representations of $Q$ over $\mathbb{C}$ are $1$-dimensional if and only if $G(Q)$ is abelian.
\end{proposition}
\begin{proof}
If $G(Q)$ is abelian all the irreducible representations of $G(Q)$ and hence of $Q$ are $1$-dimensional. If all the irreducible quandle representations of $Q$ are $1$-dimensional then all irreducible representations of $G(Q)$ are $1$-dimensional. By the previous theorem $G(Q)$ admits a faithfull unitary finite dimensional representation. In particular, this representation is faithfull and is completely reducible into a direct sum of $1$-dimensional representations. Hence, $G(Q)$ is isomorphic to a subgroup of $(\mathbb{C}^\times)^n$ for some $n>0$ and $G(Q)$ is abelian.
\end{proof}
In \cite{ab}, the authors describe finite quandles with abelian enveloping group.

\section{Irreducible representations of some quandles}
Let $n$ and $m$ be two nonzero natural numbers. Let $Q_{n,m}$ be the set $$\{x_i \vert i \in \mathbb{Z}/n\mathbb{Z} \} \cup \{y_i \vert i \in \mathbb{Z}/m\mathbb{Z} \}$$ equipped with the binary operation defined by :
$$x_i \triangleright y_j=y_{j+1},\quad y_i\triangleright x_j=x_{j+1} ,\quad x_{i}\triangleright x_j=x_j,\quad y_i\triangleright y_j=y_j.$$
With the above binary operation $Q_{n,m}$ is a quandle. These quandles appear in \cite{EM}. 

\begin{proposition}\label{ir}
Let $\rho : Q_{n,m} \to GL(V)$ be a quandle representation. $\rho$ is irreducible if and only if $\rho(x_1)$ and $\rho(y_1)$ has no common stable subsbace other than $0$ and $V$.
\end{proposition}
\begin{proof}
If $\rho(x_1)$ and $\rho(y_1)$ has no common stable subsbace other then $0$ and $V$ then the representation is clearly irreducible. Now assume that $\rho(x_1)$ and $\rho(y_1)$ has a common stable subsbace $W$ other than $0$ and $V$. We have that $\rho(x_i)=\rho(y_1)^{i-1}\rho(x_1)\rho(y_1)^{1-i}$ and $\rho(y_i)=\rho(x_1)^{i-1}\rho(y_1)\rho(x_1)^{1-i}$. Hence, $W$ is stable under any element of $\rho(Q_{n,m})$. This proves that the representation is not irreducible. Therefore, if the representation is irreducible then $\rho(x_1)$ and $\rho(y_1)$ has no common stable subsbace other than $0$ and $V$. We have proved the equivalence.
\end{proof}
\begin{proposition}\cite{GQ}
An irreducible representation of a finite quandle over $\mathbb{C}$ is finite dimensional.
\end{proposition}
\begin{proposition}\label{ir2}
Let $\rho :Q_{n,m} \to GL(V)$ be an irreducible finite dimensional quandle representation over $\mathbb{C}$. If $V$ is not $1$-dimensional then :
\item[1)]  For some primitive $d$-th root of unity $\alpha$ with some $d$ dividing $m$ and $n$ :
$$\rho(y_1)\rho(x_1)\rho(y_1)^{-1}=\alpha \rho(x_1)$$
\item[2)] $\rho(y_1)^d=\lambda id$ for some $\lambda \in \mathbb{C}^\times$.
\item[3)] If $v$ is an eigenvector of $\rho(x_1)$ associated to the eigenvalue $\beta$, then the family 
$$\{v,\rho(y_1)(v),\dots,\rho(y_1)^{d-1}(v)\}$$ form a basis of $V$ and $$\rho(x_1)(\rho(y_1)^i(v))=\frac{\beta}{\alpha^i}\rho(y_1)^i(v),$$
for $i=0,\dots,d-1$. 
\end{proposition}
\begin{proof}
As we have seen in the first section any quandle morphism $f$ induces a morphism of envelopping groups $G(f)$. As we have argued in the proof of proposition \ref{Z0}, for $x\in Q_{n,m}$ the morphism $G(L_x) : G(Q_{n,m})\to G(Q_{n,m})$ associated to the left translation by $x$ is $c_x$ the conjugacy by the generator $x$ of $G(Q_{n,m})$. Since, $L_{x_1}=L_{x_2}=L_{y_1\triangleright x_1}$, we have that $c_{x_1}=c_{y_1\triangleright x_1}$. Hence, $x_1( y_1\triangleright x_1)^{-1}$ lies in the center of $G(Q_{n,m})$. Since the representation is irreducible (and hence finite dimensional) the center of $G(Q_{n,m})$ acts by multiplication by scalars with repspect to the representation of ${G(Q_{n,m})}\to GL(V)$ induced by $\rho$ (by the universal property of $G(Q_{n,m}))$. This proves that : $$\rho(y_1)\rho(x_1)\rho(y_1)^{-1}=\alpha \rho(x_1),$$
for some $\alpha \in \mathbb{C}^\times$. We will write $y_1^n\triangleright x_1$ for the iteration $n$ times of $y_1\triangleright$ on $x_1$ : $$y_1\triangleright (y_1\triangleright( \cdots (y_1\triangleright x_1)\cdots)).$$
We have that : 
$$ \rho(y_1^n\triangleright x_1)=\rho(x_1),$$
and by the equation $\rho(y_1)\rho(x_1)\rho(y_1)^{-1}=\alpha \rho(x_1)$ we obtained before, we have that :
$$\rho(y_1^n\triangleright x_1)=\rho(y_1)^n\rho(x_1)\rho(y_1)^{-n}=\alpha^n\rho(x_1).$$
Therefore $\alpha^n=1$. This proves that $\alpha$ is an $n$-root of unity. Manipulating the equation $\rho(y_1)\rho(x_1)\rho(y_1)^{-1}=\alpha \rho(x_1)$, we get that :
$$\rho(x_1)\rho(y_1)\rho(x_1)^{-1}=\alpha^{-1} \rho(y_1).$$
Applying the same reasoning as just before with $\rho(x_1^m\triangleright y_1)$ we get that $\alpha$ is an $m$-th root of unity. This proves that $\alpha$ is a $n$-th and $m$-th root of unity. We still need to prove that $\alpha \neq 1$ to complete the proof of $1$. Assume $\alpha=1$. This implies that $\rho(x_1)$ and $\rho(y_1)$ commute and hence any eigenspace of these two automorphisms of $V$ is a common stable subspace. It follows from the proposition \ref{ir} since $\rho $ is irreducible that $\rho(x_1)$ and $\rho(y_1)$ are scalar multiple of the identity. But this contradict the same proposition since $V$ is not $1$-dimensional. Hence, $\alpha \neq 1$. This completes the proof of $1)$. We now prove $2)$. The equation in $1)$ of the proposition proves that $\rho(y_1)^d$ commutes with $\rho(x_1)$. It obviously also commute with $\rho(y_1)$. Hence, an eigenspace of $\rho(y_1)^d$ is a stable subspace for $\rho(y_1)$ and $\rho(x_1)$. It follows from proposition \ref{ir} since $\rho$ is irreducible that an eigenspace of $\rho(y_1)^d$ is equal to $V$. This proves that $\rho(y_1)^d$ is a scalar multiple of the identity. We prove $3)$. Using the equation in $1)$ one shows that if $w$ is associated to the eigenvalue $\theta$ with repsepect to $\rho(x_1)$ then $\rho(y_1)w$ is an eigenvector associated to the eigenvalue $\frac{\theta}{\alpha}$. Applying, this result and an induction one gets the equation : 
$$\rho(x_1)(\rho(y_1)^i(v))=\frac{\beta}{\alpha^i}\rho(y_1)^i(v),$$
of $3)$. Since $\alpha$ is a primitive root it follows that the vectors of the family $$\{v,\rho(y_1)(v),\dots,\rho(y_1)^{d-1}(v)\}$$ are associated to different eigenvalues of $\rho(x_1)$ and hence the family is free. Denote by $W$ the subspace spanned by this family. $W$ is stable under $\rho(x_1)$ since it consist of eigenvectors of $\rho(x_1)$. It follows from $2)$ of this proposition that $W$ is stable by $\rho(y_1)$. Hence, $W$ is stable by $\rho(x_1)$ and $\rho(y_1)$ and therefore by the proposition \ref{ir} $W=V$ since $\rho$ is irreducible. This proves that the family in $3)$ of this proposition generates $V$. We have already seen that it is free. This completes the proof of $3)$.
\end{proof}
\begin{proposition}
Let $d$ be a positive integer greater with $d>1$ divising $n$ and $m$, $\alpha$ be a primitive $d$-th root of the unity and $\lambda, \beta \in \mathbb{C}^\times$ :
\item[1)] The exists a unique quandle representation $\rho_{\alpha,\lambda,\beta} : Q_{n,m} \to GL_d(\mathbb{C})$, such that :
 $$\rho_{\alpha,\lambda,\beta}(x_1)=\begin{pmatrix}
    \beta & 0   & \dots& \dots &  0 \\
     0 &\frac{\beta}{\alpha}   & \ddots&  &\vdots \\
    \vdots & \ddots &  \ddots & \ddots &\vdots \\
\vdots &   & \ddots  &\ddots &0\\
   0 &   \dots & \dots  & 0&\frac{\beta}{\alpha^{d-1}}
\end{pmatrix} \quad   \rho_{\alpha,\lambda,\beta}(x_2)=\begin{pmatrix}
   0 & 0   & \dots& 0 &  \lambda \\
     1 & \ddots   & \ddots&  &0\\
    0 & \ddots &  \ddots & \ddots &\vdots \\
\vdots &  \ddots & \ddots  &\ddots &0\\
   0 &   \dots & 0 & 1&0
\end{pmatrix}.$$
\item[2)] The quandle representation $\rho_{\alpha,\lambda,\beta}$ is given by the following : 
$$\rho _{\alpha,\lambda,\beta}(x_i)=\alpha^{i-1}\rho_{\alpha,\lambda,\beta}(x_1) \quad \text{and} \quad \rho_{\alpha,\lambda,\beta}(y_j)=\alpha^{1-j}\rho_{\alpha,\lambda,\beta}(y_1),$$
for $i\in \mathbb{Z}/n\mathbb{Z}$ and $j\in \mathbb{Z}/m\mathbb{Z}$.
\item[3)] The commutator of $(\rho(y_1),\rho(x_1))$ is equal to $\alpha id$.
\item[4)] The representation $\rho_{\alpha,\lambda,\beta}$ is irreducible.
\end{proposition}
\begin{proof}
Let $A$ be the matrix corresponding to $\rho_{\alpha,\lambda,\beta}(x_1)$ and $B$ the one corresponding to $\rho_{\alpha,\lambda,\beta}(x_1)$ in $1)$. One checks that $BA=\alpha AB$ this implies that $AB=\frac{1}{\alpha}BA$. These equations with the defining binary operation of $Q_{n,m}$ allow to show that $\rho_{\alpha,\lambda,\beta}$ defined by equations of $1)$ and $2)$ is indeed a quandle representation. This proves the existence of $\rho_{\alpha,\lambda,\beta}$ and that $2)$ follows from the uniqueness in $1)$. The uniqueness in $1)$ follows from the fact that 
$$\rho_{\alpha,\lambda,\beta}(x_i)=\rho_{\alpha,\lambda,\beta}(y_1^{i-1}\triangleright x_1)=\rho_{\alpha,\lambda,\beta}(y_1)^{i-1}\rho_{\alpha,\lambda,\beta}(x_1)\rho_{\alpha,\lambda,\beta}(y_1)^{1-i},$$
and $$\rho_{\alpha,\lambda,\beta}(y_j)=\rho_{\alpha,\lambda,\beta}(x_1^{j-1}\triangleright y_1)=\rho_{\alpha,\lambda,\beta}(x_1)^{j-1}\rho_{\alpha,\lambda,\beta}(y_1)\rho_{\alpha,\lambda,\beta}(x_1)^{1-j},$$
where we denote the $n$-times iteration of $x\triangleright$ on $y$ by $x^n\triangleright y$. We have proved $1)$ and $2)$. $3)$ follows from the equation $BA=\alpha AB$ we obtained earlier. Now let $W$ be a subspace stable under $\rho(x_1)$ and $\rho(x_2)$. If $W\neq 0$, then $W$ contains an eigenvector $v$ of $\rho(x_1)$. But as one checks the family $\{v,\rho_{\alpha,\lambda,\beta}(y_1)(v),\dots,\rho_{\alpha,\lambda,\beta}(y_2)^{d-1}(v) \}$ spans $V$. Hence, $W=V$ if $W\neq 0$. It follows from proposition \ref{ir}, that $\rho_{\alpha,\lambda,\beta}$ is irreducible. We have proved $4)$ and completed the proof of the proposition.
\end{proof}
\begin{theorem}
\item[1)] The irreducible quandle representations of $Q_{n,m}$ of dimension greater then $1$ are up to equivalence the representations $\rho_{\alpha,\lambda,\beta}$ as in the previous proposition. 
\item[2)] Two representations $\rho_{\alpha,\lambda,\beta}$ and $\rho_{\alpha',\lambda',\beta'}$ are equivalent if and only if $\alpha=\alpha'$, $\lambda = \lambda'$ and  $\beta=\beta'\alpha^i$ for some $i$.
\end{theorem}
 \begin{proof}
$1)$ follows from $1)$ and $4)$ of the previous proposition and $2)$ and $3)$ of proposition \ref{ir2}. Assume $\rho_{\alpha,\lambda,\beta}$ and $\rho_{\alpha',\lambda',\beta'}$ are equivalent. It follows from $3)$ of the previous proposition that $\alpha=\alpha'$. $\lambda$ and $\lambda'$ are determined by the determinants of $\rho_{\alpha,\lambda,\beta}(y_1)$ and $\rho_{\alpha',\lambda',\beta'}(y_1)$. Hence, $\lambda=\lambda'$. Now $\beta$ and $\beta'$ are elements of the spectrum of $\rho_{\alpha,\lambda,\beta}(x_1)$ and $\rho_{\alpha',\lambda',\beta'}(x_1)$. Hence, $\beta=\beta'\alpha^i$ for some $i$. To complete the proof of $2)$ we still need to show that $\rho_{\alpha,\lambda,\beta}$ and $\rho_{\alpha,\lambda,\beta\alpha^i}$ are equivalent. Let $v$ be the eigenvector of $\rho_{\alpha,\lambda,\beta}(x_1)$ associated to $\beta \alpha^i$. The representation $\rho_{\alpha,\lambda,\beta}$ taken in the basis $\{v,\rho_{\alpha,\lambda,\beta}(y_1)(v),\dots,\rho_{\alpha,\lambda,\beta}(y_1)^{d-1}(v)\}$ is exactly $\rho_{\alpha,\lambda,\beta\alpha^i}$ ($\rho_{\alpha,\lambda,\beta}(y_2)^d=\lambda$). Indeed,  by $1)$ of the previous proposition, one only need to check this for the images of $x_1$ and $y_1$ and this is straitforward. This completes the proof of the theorem.

\end{proof}
\begin{thebibliography}{test}

\bibitem[EM14]{EM}
Eisermann, Michael, \newblock Quandle Coverings and Their Galois Correspondence. \newblock Fundamenta Mathematicae, vol. 225, no. 1, pp. 103–67,  2014.
\bibitem[ME18]{rep} 
Elhamdadi, M. and Moutuou, E. kaïoum M.. \newblock Finitely stable racks and rack representations. \newblock  Communications in Algebra, 46(11), 4787–4802, 2018.
\bibitem[LM21]{ab}
Victoria Lebed and Arnaud Mortier, \newblock Abelian quandles and quandles with abelian structure group. \newblock Journal of Pure and Applied Algebra, volume 225, 2021.
\bibitem[EPZ25]{nondec}
\newblock Elhamdadi, Mohamed and Senesi, Prasad and Zappala, Emanuele, \newblock On the representation theory of dihedral and cyclic quandles. \newblock Journal of Algebra and Its Applications, 2025.
\bibitem[MM26]{GQ}
Maassarani Mohamad, \newblock Groups and quandles. \newblock arxiv, 2026.
\bibitem[MMb26]{Qrep}
Massarani Mohamad, \newblock On irreducible representations of quandles \newblock arxiv, 2026.

%\bibitem[JO82]{DJ}
%D. Joyce,\newblock A classifying invariant of knots, the knot quandle \newblock J. Pure Appl. Algebra 23, pages 37–65, 1982.
%\bibitem[Br88]{Br}
%E. Brieskorn, \newblock Automorphic sets and singularities.\newblock Contemp. Math. 78, pages 45–115, 1988.
%\bibitem[FR92]{FR}
%R. Fenn, C. Rourke, Racks and links in codimension two, J. Knot Theory Ramifications 1 (4), pages 343–406, 1992.
%\bibitem[Ka87]{Sc}
%Karpilovsky, Gregory, \newblock The Schur Multiplier.\newblock Clarendon Press, 1987.


%\bibitem[Dar22]{dar}
%Jacques Darné, \newblock Nilpotent quandles. \newblock arxiv, 2022
%\bibitem[LE24]{leb}
%Lebed Victoria, Conjugation groups and structure groups of quandles, arxiv.
%\bibitem[NDCV26]{str}
%Rodr\'{i}guez-Nieto, Jos\'{e} Gregorio and Salazar-D\'{i}az, Olga Patricia and Vallejos-Cifuentes, Ricardo Esteban and Vel\'{a}squez, Ra\'{u}l, \newblock Rack representations and connections with groups representations. \newblock, Journal of Algebra and Its Applications, 2026.
%\bibitem[LM21]{ab}
%Victoria Lebed and Arnaud Mortier, \newblock Abelian quandles and quandles with abelian structure group. \newblock Journal of Pure and Applied Algebra, volume 225, 2021.
 



\end{thebibliography}


 


 \end{document}
